\chapter{Приложения}
\label{ch:appendix}

\section{Схема установки}
Схема экспериментальной установки приведена на рис. \ref{fig:ust}. Опыт выполнялся с помощью измерительного микроскопа. На столике микроскопа помещается держатель с пластинкой чёрного стекла. На пластинке лежит исследуемая линза.

Источником света служила ртутная лампа, находящаяся в защитном кожухе. Для получения монохроматического света применялся призменный монохроматор, состоящий из конденсора $K$, коллиматора (щель $S$ и объектив $O$) и призмы прямого зрения $\Pi$. Эти устройства с помощью рейтеров располагались на оптической скамье. Свет от монохроматора попадал на опак-иллюминатор (ОИ), расположенный между окуляром и объективом микроскопа -- специальное устройство для освещения объекта при работе в отражённом свете. Внутри опак-иллюминатора находилась полупрозрачная пластинка $P$, наклоненная под углом 45$^\circ$ к оптической оси микроскопа. Свет частично отражался от этой пластинки, проходя через объектив микроскопа и попадая на исследуемый объект. Пластинка могла поворачиваться вокруг горизонтальной оси $x$, а опак-иллюминатор -- вокруг вертикальной оси.

\begin{figure}[ht]
    \centering
    \includegraphics[width=0.8\linewidth]{./img/ust.png}
    \caption{Схема установки для наблюдения колец Ньютона, состоящая из монохроматора и микроскопа. Монохроматор включает в себя конденсор $K$, щель $S$, объектив $O$ и призму прямого зрения $\Pi$. Микроскоп дополнительно содержит полупрозрачную пластинку $P$.}
    \label{fig:ust}
\end{figure}

% Столик микроскопа может перемещаться в двух взаимно перпендикулярных направлениях с помощью винтов препаратоводителя. Отсчетный крест окулярной шкалы перемещается перпендикулярно оптической оси микроскопа с помощью микрометрического винта.

Оптическая схема монохроматора позволяла получить в плоскости входного окна опак-иллюминатора достаточно хорошо разделённые линии спектра ртутной лампы. Изображение щели $S$ фокусировалось на поверхности линзы объективом микроскопа, и в том же месте находилась плоскость наблюдения микроскопа, т.е. точка источника и точка наблюдения интерференции совпадали. Картина интерференции как и в случае расположения пластинки сверху, так и в данном случае не зависила от коэффициента преломления линзы и определялась величиной зазора между нижней поверхностью линзы и стеклянной пластинкой.

\section{Метод определения окулярной шкалы микроскопа}
Для определения цены деления окулярной шкалы сверху на линзу была помещена калиброванная объектная шкала. Плавно поднимая тубус, было найдено изображение миллиметровой шкалы и после совмещено с окулярной шкалой.

Объектная шкала размером 1 мм разбита на 100 делений. Используя всё поле зрения микроскопа, были отмечены, какие из самых дальних штрихов объектной шкалы лучше всего совпали со штрихами окулярной шкалы.

Оказалось, что одно деление окулярной шкалы используемого микроскопа с высокой точностью соответствует 0,1 мм.